\section{Entorno y Configuración (macOS / Linux)}

\subsection{Instalación de Utilidades GNU en Mac}
Para utilizar la suite completa de GNU (como la cabecera universal \texttt{<bits/stdc++.h>}) e igualar con precisión el entorno de evaluación del concurso, instala la última versión de GCC y las utilidades de núcleo de GNU mediante Homebrew:
\begin{lstlisting}[language=bash]
brew install gcc coreutils
\end{lstlisting}

\subsection{Alias de Compilación (\texttt{\textasciitilde/.zshrc})}
Agrega este alias al final de tu archivo de configuración de la terminal para compilar localmente de forma rápida. 

\textbf{Nota Crítica de Arquitectura (Apple Silicon):} En procesadores M1/M2/M3 bajo macOS, la suite \texttt{libsanitizer} dentro de GCC para ARM64 no cuenta con soporte nativo completo en Homebrew, lo que arroja el error \texttt{ld: library 'asan' not found}. Se han removido las banderas \texttt{-fsanitize=address,undefined}, pero se conservan los diagnósticos estrictos de tipos, desbordamientos lógicos y conversiones implícitas peligrosas:

\begin{lstlisting}[language=bash]
# Ajustar g++-15 a la version exacta instalada en tu equipo (ej. g++-14, g++-15)
alias c='g++-15 -Wall -Wconversion -Wfatal-errors -g -std=c++17'
\end{lstlisting}
Actualiza tu sesión corriendo en la terminal: \texttt{source \textasciitilde/.zshrc}.

También puedes simplemente agregarlo al final como: \texttt{echo "alias c='g++-15 -Wall -Wconversion -Wfatal-errors -g -std=c++17'" >> ~/.zshrc}

\subsection{Script de Verificación de Tipografía (\texttt{hash.sh})}
Crea el archivo \texttt{hash.sh} para validar la correcta transcripción de tus plantillas de código desde el cuaderno impreso. Este script compara el hash de 6 caracteres del código limpio con el hexadecimal de KACTL. 

\textbf{Nota sobre Portabilidad de Hashes:} Es indispensable incluir la bandera \texttt{-fpreprocessed} para evitar que el compilador inyecte macros internos específicos del procesador M2 o del sistema operativo macOS, lo cual corrompería el hash portable de KACTL:
\begin{lstlisting}[language=bash]
#!/bin/bash
g++-15 -E -dD -P -fpreprocessed - | tr -d '[:space:]' | gmd5sum | cut -c-6
\end{lstlisting}
Otorga los permisos de ejecución necesarios en tu terminal:
\begin{lstlisting}[language=bash]
chmod +x hash.sh
\end{lstlisting}
Úsalo de la siguiente forma para verificar fallos tipográficos: ./hash.sh < codigo.cpp

\subsection{Configuración Básica de Editor (\texttt{\textasciitilde/.vimrc})}
Habilita el auto-indentado limpio para código de C++ y el mapeo de escape veloz presionando las teclas \texttt{jk} o \texttt{kj} de forma sucesiva en modo inserción:
\begin{lstlisting}[language=bash]
set cin aw ai is ts=4 sw=4 tm=50 nu noeb bg=dark ru cul sy on
im jk <esc>
im kj <esc>
\end{lstlisting}

\subsection{Remapear Caps Lock de Forma Nativa en macOS}
Para evitar comandos de Linux incompatibles en macOS (como \texttt{xmodmap}), puedes cambiar el mapeo directamente desde la interfaz del sistema operativo:
\begin{enumerate}
    \item Ve a \textbf{Ajustes del Sistema} $\rightarrow$ \textbf{Teclado}.
    \item Haz clic en \textbf{Atajos de teclado...} y luego selecciona \textbf{Teclas de modificación}.
    \item Cambia la acción de la tecla \textbf{Bloqueo de mayúsculas} (\textit{Caps Lock}) a \textbf{Control} o \textbf{Escape} para una navegación ultra veloz en Vim.
\end{enumerate}

\subsection{Validar Binario Activo}
Si Homebrew actualiza o instala una versión superior de GCC, identifica con exactitud el sufijo del comando y las rutas ejecutando:
\begin{lstlisting}[language=bash]
brew info gcc
ls /opt/homebrew/bin/g++-*
\end{lstlisting}